\documentclass[12pt]{article}
\usepackage[margin=2cm]{geometry}
\usepackage{fourier}
\usepackage{tikz}
\usetikzlibrary{calc}
\usepackage[nolinks]{qrcode}

\newcommand{\smallgraderbox}[2]{
  \begin{tikzpicture}[scale=0.6, every node/.style={transform shape}]
    \begin{scope}[line width=1.5pt]
      \draw (0, 2.3) -- +(21.59, 0);
      \draw (2.3, 0) -- +(0, 2.3);
      \draw (4.5, 0) -- +(0, 2.3);
      \draw (19.0, 0) -- +(0, 2.3);
      \draw (4.5, 0.7) -- (19, 0.7);
      \node at (3.4, 0.7) {Page score};
    \end{scope}

    \node at (1.5, 1.5) {\qrcode[height=1cm]{#1}};
    \node at (20, 1.5) {\qrcode[height=1cm]{#1}};

    \foreach \x in {0,1,...,#2}{
      \coordinate (p) at (5.25 + 1.3*\x, 1.5);
      \node[draw, minimum size=0.4cm] at (p) {\x};
      \node[draw, minimum size=0.8cm, line width=0.5pt, style=dashed] at (p) {};
    };
  \end{tikzpicture}
}

\pagestyle{empty}
\begin{document}
\section{Automated score recording}

While we will continue our tradition of artisanally hand-grading
exams, we will completely automate the score entry process and return
exams electronically.  Here's how it will work.

\begin{enumerate}
\item Immediately after proctoring the exam, we will first sort the
  exams into piles by TA.  For each pile, we will remove the
  staples and sort into stacks consisting of a single page number.
  There will be a preprinted coversheet for each stack (``Sections ADB
  and ADC, Page 3") and each stack will be secured with a
  medium-sized binder clip.

  Note: The sorting by TA is largely just so we end up with piles of
  roughly equal sizes and is not required by the program.  The
  separating by page \textbf{speeds up grading immensely} but is not
  necessary either.

\item When grading, take the binder clip off one stack, grade all
  the pages, reattach the binder clip, and finally put the stack in
  the ``to be scanned'' box.  Note: You do \textbf{not} need to
  preserve the order of the pages in each stack.  Feel free to shuffle
  them!  The program does not care at all which order things are
  scanned in.

\item All grading should be done with a dark (e.g. blue, purple, or black) 
  felt-tipped marker.  The exams will be scanned and returned to the
  students in black and white.  

\item Each page will have the following at the bottom:

  \smallgraderbox{nmd123,0,8}{8}


\begin{enumerate}
  \item Write the total for the page in the ``page score'' box at left.  
  \item 
    \begin{minipage}[t]{10cm}
    Indicate the ``page score'' by making a cross in the
    \emph{smaller} box like this: 
    \end{minipage} \quad
    \begin{tikzpicture}[
      scale=0.5, 
      every node/.style={transform shape}, 
      line cap=round, 
      baseline=8]
      \node[draw, minimum size=1cm] at (0, 0) {};
      \node[draw, minimum size=2cm, line width=0.5pt, style=dashed] at (0, 0) {};
      \draw[line width=4pt] (-0.4, -0.4) .. controls (-0.06, 0.05) .. (0.4, 0.4);
      \draw[line width=4pt] (-0.4, 0.4) .. controls (0.06, 0.07) .. (0.47, -0.4);
    \end{tikzpicture}
    

  \item 
    \begin{minipage}[t]{10cm}
      If you need to change the grade, leave your initial mark as is,
      and fill the \emph{smaller} box for the correct grade like this:
    \end{minipage} \quad
    \begin{tikzpicture}[
      scale=0.5, 
      every node/.style={transform shape}, 
      line cap=round, 
      baseline=12]
      \node[draw, fill, minimum size=1cm] at (0, 0) {};
      \node[draw, minimum size=2cm, line width=0.5pt, style=dashed] at (0, 0) {};
    \end{tikzpicture}
    
  \item
    \begin{minipage}[t]{10cm}
       If you need to change the grade a \emph{second} time, fill in the
       \emph{large} box of the now desired grade like this:
    \end{minipage} \quad
    \begin{tikzpicture}[
      scale=0.5, 
      every node/.style={transform shape}, 
      line cap=round, 
      baseline=12]
      \node[draw, minimum size=1cm] at (0, 0) {};
      \node[draw, fill, minimum size=2cm, line width=0.5pt, style=dashed] at (0, 0) {};
    \end{tikzpicture}

  \item If you need to change the grade a third time, give up and
    hand the page over to Nathan for hand-entry of the score. 

\end{enumerate}

\end{enumerate}


\section{Supplies}

\begin{itemize}

\item Various tools for removing staples, including scissors.  For the
  midterms which were 6-7 pages, all methods for removing staples
  worked fine.  For the 12 page final, only cutting off the corners
  with scissors was effective.

\item 200 medium-sized binder clips.

\item 30 dark felt-tipped markers.

\end{itemize}

\end{document}
